\documentclass[10pt]{article}
\usepackage{amsmath,amssymb,amsthm,geometry,hyperref}
\geometry{margin=1in}
\newtheorem{proposition}{Proposition}
\newcommand{\R}{\mathbb R}
\title{Common ascent must be defined in the executable update geometry}
\author{RefCloth mathematical exploration, 2026-10-31}
\date{}
\begin{document}
\maketitle

\paragraph{Status and novelty boundary.}
This is a mathematically explicit, testable optimization candidate, not a
validated new attack. Multiobjective steepest descent, MGDA, PCGrad, CAGrad,
and maximin universal attacks are established. The proposed question is whether
the sign-and-clip execution geometry of a shared physical texture destroys
cross-state ascent, and whether respecting that geometry improves actual
native-exported updates. The linear program and its dual are standard convex
analysis. None of the results below proves closed-loop attack success.

\paragraph{Problem.}
Let $J_i(\theta)$ be the source objective for task--state group $i$, including
the same texture regularizer in every group. An example matching the current
code is $J_i=10D_i-0.1A-0.1V$. Let $g_i=\nabla J_i(\theta)$, with optional
fixed positive normalization $q_i=g_i/s_i$, $s_i=\max(\|g_i\|_1,\varepsilon)$.
Keep zero rows: a zero-gradient group prevents strict first-order common ascent.
The actual feasible parameter increment is the box
\[
 B(\theta,\eta)=\{\delta:\ell_j\leq\delta_j\leq u_j\},\qquad
 \ell_j=\max(-\eta,-5-\theta_j),\quad
 u_j=\min(\eta,5-\theta_j).
\]
It contains $0$. Solve the executable maximin problem
\begin{equation}\label{eq:primal}
 \gamma^*=\max_{\delta\in B}\min_i(q_i^\top\delta+c_i)
 =\max_{\delta,t}\{t:q_i^\top\delta+c_i\geq t\ \forall i,\ \delta\in B\}.
\end{equation}
Default $c_i=0$; fixed uncertainty penalties are discussed below. The output is
the primal increment $\delta^*$, not the sign of a weighted gradient.

\begin{proposition}[Sign can destroy common ascent]
Even for two linear objectives, averaging or applying a sign after PCGrad can
decrease one objective despite existence of a strict common-ascent box step.
\end{proposition}
\begin{proof}
Take $g_1=(1,0)$, $g_2=(-2,1)$ and $B=[-1,1]^2$.
The sign of their mean is $(-1,1)$, yielding gains $(-1,3)$.
Projecting each row against the original conflicting other row as in PCGrad
gives $\widetilde g_1=(1/5,2/5)$ and $\widetilde g_2=(0,1)$.
Their sum $(1/5,7/5)$ has gains $(1/5,1)$, both positive. Its sign $(1,1)$,
however, yields gains $(1,-1)$. In contrast, $\delta^*=(1/3,1)$ gives gains
$(1/3,1/3)$. It is maximin optimal: to obtain gain $t$, one needs
$\delta_1\geq t$ and $\delta_2\geq t+2\delta_1\geq3t$; since
$\delta_2\leq1$, $t\leq1/3$.
\end{proof}

\paragraph{Normalization is a separate factor.}
For the same example, $\ell_1$-normalized gradients are
$(1,0)$ and $(-2/3,1/3)$. Their mean has positive coordinates, so its sign is
$(1,1)$ and the second normalized gain is $-1/3$. The normalized maximin
solution is $(1/5,1)$ with common gain $1/5$. Thus normalization alone does
not fix execution geometry, and a normalized-mean control is necessary.

\begin{proposition}[Dual bound and common-ascent diagnostic]\label{prop:dual}
For the simplex $\Delta_n$, the optimal value in \eqref{eq:primal} is
\[
 \gamma^*=\min_{w\in\Delta_n}\left[
 \sum_j\max\{\ell_j h_j,u_j h_j\}+w^\top c\right],\qquad
 h=\sum_iw_iq_i.
\]
Any feasible primal increment supplies a lower bound; any simplex weight vector
supplies an upper bound. For a symmetric cube and $c=0$ this becomes
$\gamma^*=\eta\min_{w\in\Delta_n}\|\sum_iw_iq_i\|_1$.
With $c=0$, $\gamma^*>0$ if and only if a feasible strict common-ascent increment
exists. In the symmetric full-dimensional cube, $\gamma^*=0$ if and only if
$0\in\operatorname{conv}\{q_i\}$.
\end{proposition}
\begin{proof}
Write the finite minimum as $\min_{w\in\Delta_n}\sum_iw_i(q_i^\top\delta+c_i)$.
The two feasible sets are convex compact and the expression is affine in each
argument, so minimax interchange applies. Maximizing $h^\top\delta$ over a
box separates coordinate-wise and gives the displayed support function.
The remaining claims follow directly from feasibility, compactness, and the
fact that a norm is zero exactly at zero.
\end{proof}

\paragraph{Do not sign the optimal dual combination either.}
In the unnormalized example, optimal dual weights are $(2/3,1/3)$ and their
combination is $(0,1/3)$. Its sign $(0,1)$ yields worst gain $0$, whereas
the primal optimum has worst gain $1/3$. The fractional first coordinate is
chosen in a zero dual-gradient direction to satisfy the active task constraints.

\begin{proposition}[Few fractional coordinates suffice]
Suppose all $d$ box intervals have nonzero width. There exists an optimal basic
solution of \eqref{eq:primal} with at most $n-1$ increment coordinates strictly
between their two bounds. More sharply, if it has $k$ active task constraints,
the number of such coordinates is at most $k-1$.
\end{proposition}
\begin{proof}
The LP has $d+1$ variables $(\delta,t)$. If $p$ coordinates are strictly inside
their intervals, at most $d-p$ independent active box constraints are available.
At a basic solution, active task and box constraints together must span
$d+1$ dimensions, so $d-p+k\geq d+1$ and $p\leq k-1\leq n-1$.
An optimal basic solution exists because the box is bounded and the optimal
epigraph level is finite. This statement concerns an optimal basic solution,
not every point of an optimal face or an averaged no-regret iterate.
\end{proof}

\paragraph{Implication for four groups.}
A basic LP solution needs at most three fractional coordinates among the
thousands of cloth parameters; the others lie at executable box endpoints.
This avoids explaining a gain merely by shrinking the entire update. Report
the observed fractional count and primal/dual gap rather than assuming the
numerical solver returned the required basic solution.

\begin{proposition}[Conditional native-export lower bound]
Let $F_i(\theta)$ be a continuously relaxed native objective, and suppose on the
relevant box segment that
\[
 \|\nabla F_i(\theta)-g_i\|_1\leq\rho_i,\qquad
 \|\nabla F_i(a)-\nabla F_i(b)\|_1\leq L_i\|a-b\|_\infty.
\]
Suppose also that the native objective as a function of the actual texture
$\mathcal F_i(T)$ is $\kappa_i$-Lipschitz in texture $\ell_\infty$ distance and
the export quantizer has error at most $q$. For $\|\delta\|_\infty\leq\eta$,
the actually exported native gain is bounded below by
\[
 g_i^\top\delta-\rho_i\eta-\tfrac12 L_i\eta^2-2\kappa_iq.
\]
Consequently, setting $c_i=-\rho_i\eta-\tfrac12L_i\eta^2-2\kappa_iq$ in the
LP yields a common native gain of at least $\gamma^*$, under these assumptions.
For normalized objectives, divide each function, gradient and bound by the
same fixed positive $s_i$.
\end{proposition}
\begin{proof}
The gradient-error term follows from H\"older's inequality with the
$\ell_1/\ell_\infty$ dual pair. Integrating the gradient-Lipschitz remainder
along the segment gives $L_i\|\delta\|_\infty^2/2$. Comparing the relaxed and
quantized textures at both endpoints costs at most $2\kappa_iq$.
\end{proof}

\paragraph{Limits of the proposition.}
Finite calibration probes do not establish a global $\ell_1$ gradient-error or
curvature bound. Empirical estimates therefore do not constitute a certificate.
For 8-bit nearest rounding $q\leq1/(2\cdot255)$; truncation instead gives a
larger bound up to $1/255$, so inspect the real export implementation. If both
endpoints export the identical PNG, native gain is exactly zero even when the
relaxed linear prediction is positive. Discrete quantized deployment is not
differentiable in texture parameters; the relaxed theorem and actual PNG test
must be kept distinct. No bound above establishes behavior along future
closed-loop states or black-box target policies.

\paragraph{Primary-source prior work.}
\begin{itemize}
\item Fliege and Svaiter (2000), multiobjective descent:
\url{https://w3.impa.br/~benar/art/FliegeSvaiter2000.pdf}.
\item Sener and Koltun (2018), MGDA in multi-task learning:
\url{https://arxiv.org/abs/1810.04650}.
\item Yu et al. (2020), PCGrad:
\url{https://arxiv.org/abs/2001.06782}.
\item Liu et al. (2021), CAGrad:
\url{https://arxiv.org/abs/2110.14048}.
\item Zhang et al. (2025), dynamic maximin universal perturbation:
\url{https://arxiv.org/abs/2503.12793}.
\end{itemize}
\end{document}
